\section{The overconvergent \texorpdfstring{$\eta$}{eta}-formula}
\label{sec:eta-formula}

This section records a proposed $p$-adic identity at additive odd primes,
the overconvergent $\eta$-formula.  It plays no role in the main theorem.
We state it only as a conditional input, and we explain why the published
results one might hope to derive it from do not apply in the form needed.

\begin{definition}[$\eta$-certificate]\label{def:closure:eta-published-imports}
Let $E/\Q$ and let $p$ be an odd prime at which $E$ has additive reduction.
An \emph{$\eta$-certificate} at $(E,p)$ consists of: a nonzero $\beta$
(the proposed $U_p$-eigenvalue of an additive $p^2$-stabilization) with
$\beta^2=-p$; a leading term $T_{\mathrm{lead}}(E,p)$ of an overconvergent
expansion, with its order $\mathrm{ord}_T(E,p)\in\Z$; a value
$\etap(E)$; nonzero scalars standing for $p$ and $\#E(\Q)_{\tors}$; and
the identity
\[
  \etap(E)
  =
  \bigl(p\,\#E(\Q)_{\tors}\bigr)^{\mathrm{ord}_T(E,p)-1}
  \,\beta^{-2}\,
  T_{\mathrm{lead}}(E,p).
\]
The certificate also lists, as propositions, the background results it
is modelled on: Bella\"iche's critical-slope theory and moment
lemma~\citep{Bellaiche2012}, overconvergent modular
symbols~\citep{PollackStevens2011,ColemanMazur1998}, critical $p$-adic
$L$-functions~\citep{BenoisBuyukboduk2024}, the Eisenstein ideal at
prime-square level~\citep{LangWake2025}, component groups of N\'eron
models~\citep{Lorenzini1993,Ling1997}, and a nonvanishing and primitivity
statement for a cuspidal evaluation, which is a conjecture.
\end{definition}

\begin{theorem}[Conditional $\eta$-formula]\label{thm:closure:oc-eta-formula}
If $(E,p)$ carries an $\eta$-certificate, then
\[
  \etap(E)
  =
  \bigl(p\,\#E(\Q)_{\tors}\bigr)^{\mathrm{ord}_T(E,p)-1}
  \,\beta^{-2}\,
  T_{\mathrm{lead}}(E,p).
\]%
\footnote{The formalization records this as a conditional schema: the
identity is a field of the certificate and the conclusion is its
projection.}
\end{theorem}

\begin{proof}
The identity is part of the certificate.
\end{proof}

We do not claim that the listed background results, together with the
conjecture, imply the identity.  Two elementary obstructions show that at
least some of them do not apply as stated.

First, the relation $\beta^2=-p$ fixes the slope of $\beta$.  With the
valuation normalized by $v_p(p)=1$, it gives $2v_p(\beta)=1$, so
$v_p(\beta)=1/2$.  In weight $2$ the critical slope is $1$; this is the
normalization of Benois--B\"uy\"ukboduk~\citep{BenoisBuyukboduk2024}, for
whom the critical slope in weight $k$ is $k-1$.  So $\beta$ is not of
critical slope, and results about critical-slope eigenvalues do not apply
to it without a further comparison.

Second, the theorem of Lang--Wake~\citep[Theorem~1.1]{LangWake2025}
concerns a level prime $N$ and a residual prime $\ell$, both at least $5$,
with $\ell\mid N+1$; these numerical conditions are necessary, though not
sufficient, for an application.  They exclude the three curves at which
the formula has been proposed, whichever residual prime is chosen.  For
$36a1$ at $p=3$, the level prime would be $3<5$ (and the full level $36$
is not prime).  For $49a1$ at $p=7$, $N+1=8$ has no prime divisor at least
$5$.  For $121b1$ at $p=11$, $N+1=12$ has no prime divisor at least $5$.

Numerical values of $\etap$ at these three curves have been recorded,
but without complete input data or a reproducible computation of the
right-hand side.  We do not regard them as evidence for the formula, at
those curves or in general.%
\footnote{Lean proves the slope computation (from the relation
$\beta^2=-p$, the multiplicativity of the valuation at $\beta^2$ and
$v_p(-p)=1$, it derives $v_p(\beta)=1/2\neq1$) and the exclusion of the
level primes $3$, $7$ and $11$ from the numerical conditions of
Lang--Wake for every residual prime.}
